\chapter{Reading Instruction Descriptions}
\label{ch:reading-instructions}

This part of the manual describes each SIRCIS instruction in a consistent format. The goal is to make each instruction
entry usable as an implementation contract, not only as programmer guidance.

\section{Instruction Entry Fields}

Each instruction entry begins with a title naming the mnemonic or mnemonic family, followed immediately by one
applicability line: \texttt{SIRC-1, all revisions} for CPU instructions, or the coprocessor and revision required for
coprocessor instructions, per the model and revision policy in Chapter~\ref{ch:coprocessor}. The fields below then
follow, in this order, for every entry:

\begin{description}
	\item[Mnemonic] The assembly mnemonic or mnemonic family, given in the entry title. Immediate and register variants are
	      listed together when they perform the same operation.
	\item[Opcodes] The 6-bit opcode values for each supported instruction format. Meta-instructions that assemble to another
	      instruction's encoding use \textbf{Assembles to} instead of \textbf{Opcodes}.
	\item[Instruction Fields] One line per encoding field the instruction uses: register fields, the condition field, the status
	      update source field and which of its encodings are legal for this instruction, shift type and count, and immediate
	      width. Values come from Chapter~\ref{ch:instruction-formats} and the facts digest; an entry must not invent a value
	      here.
	\item[Syntax] The accepted assembler forms. Meta-instructions are called out separately from real CPU instruction encodings.
	\item[Operands] The permitted operand kinds and addressing modes for the instruction. Addressing-mode names come from the
	      closed set in Table~\ref{tab:addr-modes} (Chapter~\ref{ch:addressing-modes}): Immediate, Register Direct, Address
	      Register Direct, Indirect Immediate, Indirect Register, Post-Increment, and Pre-Decrement. An entry must not introduce
	      a mode name outside that set; the 8-bit short-immediate encoding is a width variant of Immediate, not a mode.
	\item[Operation] Pseudocode for the architectural effect of the instruction. Arithmetic is performed on 16-bit word values
	      unless otherwise stated.
	\item[Description] Prose description of the instruction's behavior, supplementing the Operation pseudocode.
	\item[Flags] The effect on the condition flags in the status register.
	\item[Write-back] Whether the computed result is written to a register, discarded, or used to update control state.
	\item[Exceptions] Faults or traps that can be raised by the instruction. When an exception can occur, the entry must state
	      whether partial side effects are possible.
	\item[Condition field] The condition field values the instruction accepts for predicated execution. Every normal instruction
	      may be conditional. If the condition is false, the instruction has no architectural side effects unless a later
	      instruction description explicitly states otherwise.
	\item[Timing] The base execution time and any circumstances that add wait states.
	\item[Privilege] Whether the instruction is available in protected mode, supervisor mode, or both.
	\item[Example / Examples] One or more worked examples of the instruction in use. Each example is numbered \texttt{Example
		      N-M:}, where N is the chapter number (or appendix letter) and M counts from 1 within the chapter; multiple examples in
	      one entry take consecutive numbers.
	\item[Notes] Additional informative remarks. A Notes block is informative and must not introduce a requirement that is not
	      already stated in a normative field.
\end{description}

\section{Notation}

\begin{table}[H]
	\centering
	\begin{tabular}{lp{10cm}}
		\toprule
		\textbf{Notation}          & \textbf{Meaning}                                                                         \\
		\midrule
		\texttt{rD}                & Destination general-purpose register                                                     \\
		\texttt{rS}                & Source general-purpose register                                                          \\
		\texttt{rS1}               & First source register, usually the left-hand operand                                     \\
		\texttt{rS2}               & Second source register, usually the right-hand operand                                   \\
		\texttt{rO}                & Offset or index general-purpose register supplying a register displacement in a
		memory or effective-address form                                                                                      \\
		\texttt{\#imm16}           & 16-bit immediate word encoded in immediate format                                        \\
		\texttt{\#imm8}            & 8-bit immediate encoded in short immediate format                                        \\
		\texttt{shift}             & Optional shift definition. It applies to the register operand being shifted (the first
		source register in register and short-immediate formats, or the offset register for register-displacement
		\mnemonic{LOAD}/\mnemonic{STOR} forms); it never applies to an immediate value, a second source register, or
		loaded data. See Chapter~\ref{ch:shift-operations} for the full encoding.                                             \\
		\texttt{SR.X}              & Status register field \texttt{X}                                                         \\
		\texttt{addr / src}        & Source address-register pair (\reg{l}, \reg{a}, \reg{s}, or \reg{p}) supplying the base
		address for a memory-class or address-calculation instruction                                                         \\
		\texttt{dest}              & Destination address-register pair for an address-calculation instruction, encoded in the
		register field                                                                                                        \\
		\texttt{\#offset / \#disp} & Displacement added to the low word of an address-register pair                           \\
		\texttt{@label}            & Symbol reference resolved by the linker; PC-relative for \mnemonic{BRAN} and
		\mnemonic{BRSR}, an absolute low-word address otherwise                                                               \\
		\texttt{|cond}             & Condition-field suffix selecting the predicate for conditional execution                 \\
		\bottomrule
	\end{tabular}
	\caption{Instruction Description Notation}
	\label{tab:instruction-description-notation}
\end{table}

\section{Status Flag Effect Symbols}
\label{sec:flag-effect-symbols}

Flag tables use the following symbols:

\begin{table}[H]
	\centering
	\begin{tabular}{cp{9cm}}
		\toprule
		\textbf{Symbol} & \textbf{Meaning}                                                                      \\
		\midrule
		\texttt{*}      & Updated from the instruction result                                                   \\
		\texttt{0}      & Cleared to zero                                                                       \\
		\texttt{1}      & Set to one                                                                            \\
		\texttt{-}      & Preserved                                                                             \\
		\texttt{S}      & Updated from the shifter result when \texttt{[S]} is used                             \\
		\texttt{U}      & Architecturally undefined; the implementation leaves the flag in an unspecified state \\
		\bottomrule
	\end{tabular}
	\caption{Status Flag Effect Symbols}
	\label{tab:flag-effect-symbols}
\end{table}

\section{Conditional Execution}

Every encoded instruction contains a condition-code field. If the condition evaluates true, the instruction executes
normally. If the condition evaluates false, the instruction is skipped: destination registers, address registers,
memory, the program counter, link register, coprocessor state, and status flags are not modified by that instruction.

\section{Status Update Overrides}

ALU instructions can select the source of status flag updates with the status update field:

\begin{table}[H]
	\centering
	\begin{tabular}{lll}
		\toprule
		\textbf{Syntax} & \textbf{Encoded AF} & \textbf{Effect}                      \\
		\midrule
		\texttt{[N]}    & \texttt{00}         & Preserve status flags                \\
		\texttt{[A]}    & \texttt{01}         & Update flags from the ALU result     \\
		\texttt{[S]}    & \texttt{10}         & Update flags from the shifter result \\
		Reserved        & \texttt{11}         & Reserved for future use              \\
		\bottomrule
	\end{tabular}
	\caption{Status Update Source Encoding}
	\label{tab:status-update-source-encoding}
\end{table}

When no override is written, ALU instructions use \texttt{[A]}. The one exception is \mnemonic{SHFT}, which defaults to
\texttt{[S]}. Instructions that do not support status updates must not use status update override syntax.
